\chapter{Migrating from C to Rust}
\label{chap:ch7}

\par Gentlix was not produced by an automatic translator. The Rust backend was written 
to mirror an existing C design: same schema, same routes, same JSON. This chapter 
summarises why that manual path was necessary and what a developer moving from the C 
codebase should expect cognitively.

\section{The Cognitive Gap}
\label{sec:ch7sec1}

\par C programmers think in terms of pointers, explicit lifetimes, and conventions 
the compiler does not enforce. Rust asks for ownership up front: every value has one 
owner, references borrow for a limited time, and mutation through shared references 
is restricted. The first weeks on Gentlix often produced borrow-checker errors on 
seemingly reasonable code---for example, storing a reference to a row buffer that 
outlives the database result.

\par McNamara~\cite{McNamara2021} describes systems programming in Rust as learning to 
work \emph{with} the compiler rather than fighting it. In Gentlix, \textbf{ownership} 
replaced manual \texttt{malloc}/\texttt{free} in services and JSON building; 
\textbf{traits} replaced hand-wired vtables but required \texttt{async\_trait} and 
\texttt{Send} bounds on repositories; \textbf{async I/O} added \texttt{await} for 
\texttt{sqlx} and \texttt{spawn\_blocking} for Argon2 and analytics---a different 
model than libmicrohttpd's blocking threads even when the product behaviour matched.

\par Flitton and Morton~\cite{Flitton2023} explain async migration pitfalls---blocking 
\texttt{libpq} on Tokio tasks, or mixing \texttt{Send} bounds with legacy pointer APIs. 
Drysdale~\cite{Drysdale2024} recommends learning \texttt{Result} and iterator patterns 
before porting large C modules; Gentlix ported service traits first, then HTTP, then filled 
in domain edge cases to match C JSON field-for-field.

\par Palmieri~\cite{Palmieri2022} recommends introducing Rust incrementally in backend 
teams. Gentlix effectively did the extreme version: a full rewrite of one backend 
while keeping the C version as a reference implementation. That made comparisons fair 
but doubled maintenance during development.

\section{Automated Migration}
\label{sec:ch7sec2}

\par Tools such as \texttt{c2rust} can translate C syntax into Rust mechanically, but 
the output is usually \texttt{unsafe}, idiomatically foreign, and still missing domain 
structure. Gentlix relies on opaque pointers, vtables, and consistent error enums---patterns 
Schreiner~\cite{Schreiner1993} documents well in C---none of which become safe Rust 
modules without human redesign.

\par Automated translation also cannot recover financial invariants. Cross-backend JWT 
compatibility, integer-cent ledger rules, SQL transaction boundaries, and byte-for-byte 
API parity required deliberate choices in both languages. A migrated skeleton might compile 
after heavy \texttt{unsafe}, but it would not satisfy the non-functional requirements in 
Chapter~\ref{sec:ch2sec1sub2} without the same service-layer work already present in 
Gentlix.

\par A realistic migration path for a similar project would look like Gentlix in reverse: 
keep PostgreSQL and HTTP contracts fixed, port the domain and repository traits first, 
wrap existing C only where needed, and replace hot paths after tests prove parity. Full 
automation is useful for spotting structural gaps, not for shipping a banking API unchanged.

\section*{Summary}

\par Moving from Gentlix C to Gentlix Rust is feasible because the architecture was 
layered from the start, but it is not free. The cognitive gap shows up in ownership, 
async, and trait bounds; automated tools do not remove the need to re-prove safety and 
API parity. The dual-backend repository is therefore both a comparison instrument and 
a record of what a disciplined manual migration entails.
